% !TeX root = ../complete.tex
\section{شکل موج، معیارهای زمانی و پاسخ گذرا}
\label{sec:delay-transient}

در \NotesTopic{sec:dc-transfer}{مشخصه انتقال \lr{DC}}، خروجی ماندگار
را برای هر ولتاژ ورودی تعیین کردیم. در تحلیل گذرا، زمان رسیدن به سطح
منطقی معتبر نیز اهمیت دارد. گیتی که گره خروجی را شارژ یا تخلیه می‌کند،
\lr{Driver} است. گیت‌های متصل به آن گره و سیم رابط، بخشی از \lr{Load}
آن هستند. ظرفیت‌های خود \lr{Driver} نیز در تعیین جریان لازم برای تغییر
ولتاژ خروجی سهم دارند \cite[بخش‌های \secnum{۴.۱.۱} و
\secnum{۴.۲}، صفحات ۱۴۱ تا ۱۴۴]{weste}.

\subsection{اندازه‌گیری زمان روی شکل موج}
\label{sec:delay-waveform}

برای شکل موجی که میان زمین و \(V_{DD}\) تغییر می‌کند، مرجع معمول
اندازه‌گیری تأخیر، \(V_{DD}/2\) در هر دو شکل موج است. فاصله زمانی
بین عبور ورودی و عبور خروجی از این مرجع، تأخیر همان گذار است.
برای شرایط مشخص، بیشترینِ تأخیرهای معتبر را \lr{Propagation Delay}
با نماد \(t_{\mathrm{pd}}\) و کمترین را \lr{Contamination Delay}
با نماد \(t_{\mathrm{cd}}\) می‌نامیم. این دو حد، دیرترین و زودترین
عبور خروجی از مرجع را مشخص می‌کنند. در برآورد دقیق‌تر می‌توان حدها
را برای هر ورودی و برای گذار صعودی یا نزولی خروجی جداگانه مشخص کرد.
برای نمونه، \(t_{\mathrm{pdf}}\) حد بالای تأخیر گذار نزولی خروجی است
\cite[بخش \secnum{۴.۱.۱} و شکل \secnum{۴.۱}، صفحات ۱۴۱ و ۱۴۲]{weste}.

\lr{Rise Time} با نماد \(t_{\mathrm{r}}\) زمان عبور یک شکل موج از
۲۰ تا ۸۰ درصد دامنه تغییر و \lr{Fall Time} با نماد \(t_{\mathrm{f}}\)
زمان عبور آن از ۸۰ تا ۲۰ درصد است. این بازه‌ها قرارداد کتاب‌اند.
میانگین \((t_{\mathrm{r}}+t_{\mathrm{f}})/2\) را کتاب به‌عنوان
\lr{Edge Rate} با واحد زمان معرفی می‌کند. در شکل
\ref{fig:delay-waveform}، پیکان بالا فاصله میان دو شکل موج و پیکان
پایین فاصله میان دو نقطه از همان شکل موج خروجی را نشان می‌دهد
\cite[بخش \secnum{۴.۱.۱}، صفحه ۱۴۱]{weste}.

\begin{figure}[htbp]
  \centering
  \begin{latin}
  \begin{tikzpicture}[x=1.52cm,y=2.35cm]
    \draw[->] (0,0) -- (7,0) node[right] {\(t\)};
    \draw[->] (0,0) -- (0,1.25) node[above] {\(V/V_{DD}\)};
    \foreach \level/\ticklabel in {0.2/۰٫۲,0.5/۰٫۵,0.8/۰٫۸,1/۱} {
      \draw[gray!60,thin,dotted] (0,\level) -- (6.65,\level);
      \node[left] at (0,\level) {\fontspec[Path=C:/Windows/Fonts/]{arial.ttf}\ticklabel};
    }
    \draw[very thick] (0,0) -- (0.5,0) -- (2.5,1) -- (6.65,1);
    \draw[very thick,blue!70!black] (0,1) -- (3,1) -- (6,0) -- (6.65,0);
    \node[above] at (5.7,1) {\(V_A\)};
    \node[above,blue!70!black] at (0.85,1) {\(V_Y\)};
    \draw[densely dashed] (1.5,0.5) -- (1.5,1.22);
    \draw[densely dashed] (4.5,0.5) -- (4.5,1.22);
    \fill (1.5,0.5) circle (1.6pt) (4.5,0.5) circle (1.6pt);
    \draw[<->] (1.5,1.18) -- (4.5,1.18) node[midway,above] {\(\Delta t_{50}\)};
    \draw[densely dashed] (3.6,0.8) -- (3.6,-0.21);
    \draw[densely dashed] (5.4,0.2) -- (5.4,-0.21);
    \fill[blue!70!black] (3.6,0.8) circle (1.6pt) (5.4,0.2) circle (1.6pt);
    \draw[<->] (3.6,-0.17) -- (5.4,-0.17) node[midway,below] {\(t_{\mathrm{f}}\)};
  \end{tikzpicture}
  \end{latin}
  \caption{شکل موج شماتیک ورودی \(V_A\) و خروجی \(V_Y\) یک \lr{Inverter}.
  \(\Delta t_{50}\) تأخیر یک گذار و \(t_{\mathrm{f}}\) زمان افت همان خروجی است.
  قرارداد اندازه‌گیری از شکل \secnum{۴.۱} کتاب \cite{weste} پیروی می‌کند.}
  \label{fig:delay-waveform}
\end{figure}

برای مثال، اگر ورودی در \qty{20}{\pico\second} و خروجی نزولی در
\qty{55}{\pico\second} از مرجع نصف تغذیه بگذرند، تأخیر این گذار
\qty{35}{\pico\second} است. اگر خروجی در
\qty{43}{\pico\second} از ۸۰ درصد و در \qty{70}{\pico\second}
از ۲۰ درصد عبور کند، زمان افت آن \qty{27}{\pico\second} است.
برای یافتن \(t_{\mathrm{pd}}\) و \(t_{\mathrm{cd}}\) گیت، دامنه
ورودی‌ها، گذارها، بار و شرایط کاری مورد بررسی باید مشخص باشند.

\subsection{تأخیر گیت و زمان رسیدن در شبکه منطقی}
\label{sec:delay-arrival-time}

زمانی را که گذار به گره \(i\) می‌رسد با \(a_i\) نمایش می‌دهیم.
در مدل مقدماتی یک شبکه ترکیبی، \(a_i\) را با دیرترین زمان رسیدن
به ورودی‌های گیتِ راه‌انداز گره \(i\) و حد بالای تأخیر همان گیت
تخمین می‌زنیم. مجموعه این ورودی‌ها را \(\operatorname{fanin}(i)\)
و حد بالای مشترک تأخیر گیت را \(t_{\mathrm{pd},i}\) می‌نامیم
\cite[رابطه \secnum{۴.۱}، صفحه ۱۴۲]{weste}.
\begin{equation}
  a_i=\max_{j\in\operatorname{fanin}(i)} a_j+t_{\mathrm{pd},i}.
  \label{eq:delay-arrival}
\end{equation}
اگر زمان رسیدن دو ورودی \qty{25}{\pico\second} و
\qty{40}{\pico\second} و تأخیر گیت \qty{30}{\pico\second} باشد،
مدل، زمان رسیدن خروجی را \qty{70}{\pico\second} می‌دهد. اگر
خروجی تا \qty{85}{\pico\second} لازم باشد، \lr{Slack}، یعنی
زمان لازم منهای زمان رسیدن، \qty{15}{\pico\second} است. در تحلیل
دقیق‌تر، تأخیر هر ورودی، جهت گذار و امکان ایجاد گذار در خروجی
بررسی می‌شود \cite[بخش \secnum{۴.۱.۱}، صفحه ۱۴۲]{weste}.

مسیرهایی که مهلت زمانی را محدود می‌کنند \lr{Critical Path} هستند.
می‌توان تأخیر را در سطح \lr{Microarchitecture} با تغییر تقسیم کار میان
مراحل، در سطح منطق با انتخاب توپولوژی، در سطح مدار با اندازه‌گذاری
ترانزیستورها و در سطح \lr{Layout} با کاهش بار سیم و اتصال‌ها تغییر داد.
این فصل بر تعداد مراحل، نوع گیت و اندازه‌گذاری در سطح منطق و مدار
تمرکز دارد \cite[بخش \secnum{۴.۱.۲}، صفحات ۱۴۲ و ۱۴۳]{weste}.

\subsection{چرا تغییر ولتاژ خروجی زمان می‌برد؟}
\label{sec:delay-inverter-transient}

خروجی \(Y\) یک \lr{Inverter} را در نظر بگیرید که با یک سیم کوتاه،
ورودی گیت بعدی را راه‌اندازی می‌کند. در تقریب ظرفیت متمرکز، مقاومت
سیم ناچیز فرض می‌شود و ظرفیت سیم، ظرفیت‌های \lr{Diffusion} متصل به
\(Y\) و ظرفیت ورودی گیت بعدی با هم جمع می‌شوند. ظرفیت معادل حاصل را
\(C_{\mathrm{out}}\) می‌نامیم. ظرفیت یک پایانه متصل به گره دیگری
به‌صورت مستقیم در این جمع وارد نمی‌شود. برای مثال، ظرفیت
\lr{Source-to-Body} با دو پایانه در ولتاژ ثابت، در این گذار بار
اضافی از خروجی نمی‌گیرد \cite[بخش \secnum{۴.۲} و شکل
\secnum{۴.۳}، صفحات ۱۴۳ و ۱۴۴]{weste}.

ورودی را ابتدا صفر بگیرید. خروجی در \(V_{DD}\) است. پس از گذار
پله‌ای ورودی به \(V_{DD}\)، در مدل ایدئال \lr{pMOS} خاموش و
\lr{nMOS} روشن است. اگر \(I_n\) بزرگی جریان خروجی به زمین و
\(C_{\mathrm{out}}\) ثابت باشد، تعادل جریان گره می‌دهد
\begin{equation}
  C_{\mathrm{out}}\frac{\mathrm{d}V_Y}{\mathrm{d}t}=-I_n.
  \label{eq:delay-step-kcl}
\end{equation}
علامت منفی یعنی جریان خروجی، بار خازن و به‌دنبال آن \(V_Y\) را کاهش
می‌دهد. واحد هر دو طرف آمپر است. این رابطه صورت لحظه‌ای همان
\NotesTopic{eq:impl-charge-delay}{رابطه بار، جریان و زمان} است.

در مدل \NotesTopic{sec:book-mos-current}{\lr{Long-channel}}، اگر
\(V_{DD}>V_{tn}\)، اضافه‌ولتاژ \(V_{GT}=V_{DD}-V_{tn}\) است.
وقتی \(V_Y\) از \(V_{GT}\) بیشتر است، ترانزیستور در \lr{Saturation}
قرار دارد. پس از رسیدن خروجی به \(V_{GT}\)، ترانزیستور وارد
\lr{Linear Region} می‌شود. با صرف‌نظر از اثرهای غیرایده‌آل،
\(\beta_n=\mu_n C_{\mathrm{ox}}W/L\) ضریب فناوری و هندسه است و
رابطه جریان را می‌توان با شرط هر ناحیه نوشت \cite[رابطه
\secnum{۴.۴}، صفحه ۱۴۴]{weste}.
\begin{equation}
  \label{eq:delay-step-current}
  I_n=
  \begin{cases}
    \dfrac{\beta_n}{2}V_{GT}^{2},
      & V_Y\ge V_{GT},\\[3pt]
    \beta_n\left(V_{GT}V_Y-\dfrac{V_Y^2}{2}\right),
      & 0\le V_Y<V_{GT}.
  \end{cases}
\end{equation}
دو عبارت در مرز \(V_Y=V_{GT}\) برابرند. در ناحیه اول، جریان
ثابت است و از رابطه \eqref{eq:delay-step-kcl}، ولتاژ با شیب ثابت
\(-\beta_n V_{GT}^2/(2C_{\mathrm{out}})\) پایین می‌آید. با ورود به
ناحیه دوم، جریان با کاهش \(V_Y\) کم می‌شود. اگر
\(V_{DD}/2<V_{GT}\) باشد، رسیدن خروجی به مرجع نصف تغذیه پس از
ورود به ناحیه دوم رخ می‌دهد و باید سهم زمانی هر دو ناحیه را حساب کرد.
مدل \lr{RC} در بخش \ref{sec:delay-rc-model} راهی برای تخمین زمان
عبور بدون حل جداگانه معادله هر ناحیه فراهم می‌کند.

وقتی ورودی به‌صورت \lr{Ramp} بالا می‌رود، جریان هر دو ترانزیستور
به \(V_A(t)\) وابسته است. برای \(V_{tn}+|V_{tp}|<V_{DD}\)، ابتدا
\lr{nMOS} خاموش است. میان \(V_{tn}\) و \(V_{DD}-|V_{tp}|\)،
هر دو ترانزیستور هدایت می‌کنند و جریان خالص خازن \(I_p-I_n\) است.
در این عبارت، \(I_p\) بزرگی جریان از تغذیه به خروجی است. بالاتر از
این بازه، \lr{pMOS} خاموش می‌شود و تخلیه از مسیر \lr{nMOS} ادامه
می‌یابد. در نتیجه، هنگام مقایسه تأخیرها باید شکل موج ورودی و بار
خروجی را مشخص کرد \cite[بخش \secnum{۴.۲} و جدول
\secnum{۴.۱}، صفحات ۱۴۵ و ۱۴۶]{weste}.
